\subsection{A Last-Bar Same-Day Option Reading}
\label{sec:close_option_results}

The forecasts of Table~\ref{tab:block_ladder} are one-bar forecasts of
realized variance. This subsection reads them against the last thirty
minutes of the expiration-day chain, under the rules of
Section~\ref{sec:close_option_methods}: the nearest out-of-the-money
call and put taken at the 15:30~ET mid, cash-settled at the close, with
the position fixed by the sign of
$s_t = \widehat{RV}_t - \mathrm{IV}_{30,t}^{2}$. Joining the chain to
the eight forecast columns leaves 866 expiration days common to all
eight (3~January~2020 to 30~April~2024). The 413 expiration days
after that date --- the most recent twenty months of chain --- carry no
forecast row and appear nowhere below.
Table~\ref{tab:close_option_rules} crosses every rule with every forecast
on the days that remain.

\begin{table}[H]
\centering
\caption{Last-bar same-day option trade by position rule, 15:30~ET to
the close (Section~\ref{sec:close_option_methods}). One panel per rule;
rows are the eight forecast columns. One return per expiration day, on the
866 days common to all eight (3~January~2020 to 30~April~2024),
filled at the 15:30 mid. $t=\sqrt{n}\,\mathrm{mean}/\mathrm{std}$ is
computed on the daily $R'$ and is not annualized;
$\mathrm{Sharpe}_{\mathrm{ann}}=
(\mathrm{mean}/\mathrm{std})\times\sqrt{252}$ is the only annualized
column, and it scales one \emph{trade} day to a year. The 0DTE frame
trades about 200.3 days a year over this sample rather than $252$ --- SPXW
listed Monday, Wednesday and Friday expirations before June 2022 and every
session after --- so a calendar-time annualization multiplies that column
by 0.892 throughout and changes no comparison inside it.
\emph{ex.\ kurt} is the bias-corrected sample excess kurtosis
(a Gaussian scores 0). Every $t$ in this subsection is the plain
$\sqrt{n}\,\mathrm{mean}/\mathrm{std}$ of a daily series unless it is
named heteroskedasticity-robust or paired. The
always-short rule takes no forecast, so its eight model rows are
identical and it is shown as a single row. $\dagger$ marks the four
columns still fitted on the earlier design panel, for which a run on the
panel of record is pending
(Section~\ref{sec:close_option_methods}).
\label{tab:close_option_rules}}
\resizebox{\textwidth}{!}{% AUTO-GENERATED by writeup/make_rule_by_strategy_tex.py — do not edit.
% Source: results/atm_straddle_0dte_1530/rule_by_strategy_*.csv
\begingroup\small\setlength{\tabcolsep}{4.5pt}
\begin{tabular}{lrrrrrrrrr}
\toprule
& $n$ & mean & std & min & max & skew & ex.\ kurt & $t$ & Sharpe$_{\mathrm{ann}}$ \\
\midrule
\multicolumn{10}{l}{\emph{Short volatility (always short)}} \\
all models (no forecast) & 866 & 0.014 & 1.128 & -10.32 & 1.00 & -2.31 & 10.89 & 0.38 & 0.20 \\
\midrule
\multicolumn{10}{l}{\emph{$\mathrm{sign}(s)$ ($\pm 1$ on the sign of $s$)}} \\
baseline (HAR + calendar OLS) & 866 & 0.069 & 1.126 & -5.42 & 10.32 & 0.58 & 10.71 & 1.79 & 0.97 \\
block-diagonal ridge & 866 & 0.095 & 1.124 & -5.42 & 10.32 & 0.92 & 10.61 & 2.48 & 1.34 \\
block-diagonal ridge, without the FOMC columns & 866 & 0.090 & 1.124 & -5.42 & 10.32 & 1.00 & 10.59 & 2.36 & 1.27 \\
LightGBM$^{\dagger}$ & 866 & 0.103 & 1.123 & -5.42 & 10.32 & 1.01 & 10.58 & 2.71 & 1.46 \\
XGBoost$^{\dagger}$ & 866 & 0.106 & 1.123 & -5.42 & 10.32 & 0.91 & 10.62 & 2.78 & 1.50 \\
lasso (causally tuned)$^{\dagger}$ & 866 & 0.100 & 1.123 & -6.84 & 10.32 & 0.46 & 10.78 & 2.62 & 1.41 \\
lasso ($\alpha=10^{-4}$) & 866 & 0.107 & 1.123 & -5.42 & 10.32 & 1.22 & 10.50 & 2.80 & 1.51 \\
elastic net (causally tuned)$^{\dagger}$ & 866 & 0.072 & 1.125 & -6.84 & 10.32 & 0.46 & 10.74 & 1.90 & 1.02 \\
\bottomrule
\end{tabular}
\endgroup
}
\end{table}

Read the control first. Selling the package every day earns a mean
0.014 per expiration day at a standard deviation of 1.128
($t=0.38$, annualized Sharpe 0.20): a positive premium, at a size that
866 days cannot separate from zero. Putting the forecast in charge of
\emph{direction} is what pays, and it pays in the same direction in every
column: the rule $q=\mathrm{sign}(s)$ earns between 0.069 and
0.107 per day across the eight, and every one of the eight improves
on the control.

How \emph{strong} that evidence is is a different question from how
consistent it is, and the honest answer is that it is not strong. The two
rules trade the same contracts on the same days, so the statistic that
carries the comparison is the paired one --- the daily difference between
the two portfolios --- and not each portfolio's own $t$ against zero.
Paired against the control, the eight improvements run from 0.76 to
1.31 in annualized Sharpe at an autocorrelation-robust $t$ between
1.18 and 1.90 (plain paired $t$ 1.06 to
1.70), and seven of the eight bootstrap percentile intervals on
that difference include zero. The exception is XGBoost, whose interval
clears zero by a knife edge --- a lower bound of 0.017 on an
interval 2.46 wide. For the column the rest of the paper argues for
the difference is 1.13 at an autocorrelation-robust paired $t$ of
1.57. Two
statistics that would read as significance are worth naming so they are not
mistaken for it: each portfolio's own $t$ against zero clears $1.96$ for
six of the eight, and the active-return $t$ against the control
clears it for none of them. Eight columns pointing the same way is worth
more than any one of them --- though not eight times more, since they are
close to one bet (below) --- but on this sample only one of them separates
itself from the always-short control, and that marginally. What the panel
establishes is the rule, at that strength, and not any model.

The headline column is the block-diagonal ridge fitted on the design that
carries the FOMC calendar block, the panel of record of
Section~\ref{sec:close_option_methods}. It earns a mean of 0.095,
$t=2.48$, and an annualized Sharpe of 1.34 --- with a bootstrap
standard error of 0.55 and a 95\% percentile interval from
0.27 to 2.42, which is the width an 866-day Sharpe ratio
has and the reason no ordering below is asserted. Directly beneath it sits
the same estimator on the earlier design, which has no calendar block, and
the gap between those two rows is what the calendar channels are worth on
this trade: +0.07 of annualized Sharpe at a paired $t$ of 0.26, on
an interval from $-0.49$ to $0.63$ --- one-signed, small, and not
resolved. It is reported because the headline column and four of its
comparators do not sit on the same design, and a reader is entitled to see
the size of that.

The ridge does not lead the panel: the fixed lasso (1.51), XGBoost (1.50), LightGBM (1.46) and the causally tuned lasso (1.41) all score higher. The paper does not read
that ordering as a ranking of forecasts. The eight columns take the same
side of the trade on 82 to 87 percent of days, which is
far too little separation for 866 days to order them by; the spread
down the column is smaller than the paired intervals just quoted, and the
ordering is left unresolved. The forecast sits above the quoted implied
variance --- a long-volatility day --- on 32 to 41 percent of
days depending on the column.

The qualifications matter as much as the point estimates.

\paragraph{Where the level comes from.} Two sample splits qualify the
level rather than the sign. Dropping the 39 expiration days through
31~March~2020 --- the February--March 2020 spike --- \emph{raises} the
block-diagonal ridge portfolio to Sharpe 1.51 ($t=2.73$ on the
remaining 827 days): the full-sample figure is dragged, not
flattered, by that quarter. Splitting instead at
16~May~2022, when SPXW listings became daily --- a market-structure date, not a
threshold chosen on the portfolio --- the same portfolio scores Sharpe 0.98
($t=1.20$) on the 377 days before and 1.63 ($t=2.27$) on the
489 after: the direction edge is positive on both tapes but earns its
significance on the daily-0DTE era. The split is not uniform across columns
--- the two tree forecasts score higher before the date than after --- which is
one more reason to read the panel as one rule rather than as eight results.

\paragraph{Settlement pins carry a disproportionate share.} The
package expires worthless --- $R=-1$ exactly --- on 193 of the 866
days, 22.3 percent, whenever the close settles inside the one-strike
window. The always-short control collects $+1$ on those days against
$-0.27$ on all others; the $\mathrm{sign}(s)$ portfolio averages
$+0.18$ on them against $+0.07$ elsewhere, so pin days supply
43 percent of its total return from 22.3 percent of the
days. The portfolio monetizes a direction
call and the pin premium together, and its P\&L remains tilted toward
days the index settles inside the strikes.

\paragraph{Eight rows are one or two bets.} The eight daily $\mathrm{sign}(s)$
portfolios are correlated (mean pairwise correlation 0.70, from
0.58 to 0.93; their buy-day sets overlap at a mean pairwise Jaccard
index of 0.70), an effective count of 1.8 independent portfolios by the
participation ratio $(\sum_i \lambda_i)^2 / \sum_i \lambda_i^2$ of their
correlation matrix. Agreement down a column of
Table~\ref{tab:close_option_rules} is evidence that the rule is insensitive to
which forecast drives it --- not eight independent confirmations of the rule.

% Decision-aligned scoring of the 15:30 forecasts. Every number below is a
% macro written by writeup/make_murphy_1530_tex.py from the CSVs of
% writeup/intraday_proposals/59_trade_aligned_loss.py; the generator also
% asserts each qualitative claim made about the figure.
\ifdefined\murOracleAgree\else% GENERATED FILE -- do not edit by hand.
% Written by writeup/make_murphy_1530_tex.py from
% results/atm_straddle_0dte_1530/proposals/59/*.csv.
\newcommand{\murAgreePayMax}{55.3}
\newcommand{\murAgreePayMin}{54.0}
\newcommand{\murBaseLeadN}{three}
\newcommand{\murBaseLower}{1.54}
\newcommand{\murBaseRegionN}{five}
\newcommand{\murBaseUpper}{2.18}
\newcommand{\murBestSlice}{LightGBM}
\newcommand{\murBestSliceScore}{0.094}
\newcommand{\murGridMax}{4}
\newcommand{\murGridMin}{0.25}
\newcommand{\murHitVarMax}{68.8}
\newcommand{\murHitVarMin}{64.7}
\newcommand{\murLassoFixedQLIKERank}{eighth}
\newcommand{\murLassoFixedShRank}{first}
\newcommand{\murLassoFixedShX}{1.04}
\newcommand{\murLassoFixedSliceRank}{fifth}
\newcommand{\murLinUpper}{0.71}
\newcommand{\murNForecasts}{eight}
\newcommand{\murNLadder}{fifteen}
\newcommand{\murNUnion}{23}
\newcommand{\murOracleAgree}{62.9}
\newcommand{\murOracleMiss}{37.1}
\newcommand{\murOracleShMid}{4.53}
\newcommand{\murOracleShX}{4.07}
\newcommand{\murPayoffOracleShX}{17.3}
\newcommand{\murRhoES}{\ensuremath{0.40}}
\newcommand{\murRhoESHi}{\ensuremath{0.90}}
\newcommand{\murRhoESLo}{\ensuremath{-0.62}}
\newcommand{\murRhoQLIKE}{\ensuremath{0.12}}
\newcommand{\murRhoQLIKEHi}{\ensuremath{0.76}}
\newcommand{\murRhoQLIKELo}{\ensuremath{-0.69}}
\newcommand{\murRhoUnionES}{\ensuremath{0.97}}
\newcommand{\murRhoUnionQLIKE}{\ensuremath{0.95}}
\newcommand{\murRidgeNoFomcRank}{eighth}
\newcommand{\murRidgeNoFomcScore}{0.111}
\newcommand{\murRidgeQLIKERank}{sixth}
\newcommand{\murRidgeRank}{sixth}
\newcommand{\murRidgeSliceScore}{0.109}
\newcommand{\murTieLower}{2.38}
\newcommand{\murTreeLower}{0.77}
\newcommand{\murTreeUpper}{1.41}
\fi
\paragraph{What the forecasts are scored on, and what the trade uses.} The
columns of Table~\ref{tab:close_option_rules} cannot be ordered by the trade,
and the paper's accuracy measure does not order them on it either. The reason
is structural. QLIKE (Equation~\eqref{eq:dp_qlike}) is consistent for the
conditional mean of the variance \citep{Patton2011}, and every consistent
scoring function for a mean can be taken apart by threshold. For a threshold
$\theta > 0$ let
\begin{equation}
    S_\theta(f, y) \;=\; |y - \theta|\;
    \mathbf{1}\{\min(f, y) \le \theta < \max(f, y)\},
    \label{eq:elementary_score}
\end{equation}
the loss of someone who uses the forecast $f$ to decide whether the outcome
$y$ will exceed $\theta$: nothing when $f$ and $y$ fall on the same side of
$\theta$, and the distance by which $y$ landed from $\theta$ when they do not.
Every consistent scoring function for the mean is, up to a term that does not
involve the forecast, a mixture of these elementary scores over $\theta$ with
a nonnegative weight \citep{Ehm2016}; for a Bregman loss generated by a convex
function $\phi$ the weight is $\phi''(\theta)\,d\theta$ in the normalization of
Equation~\eqref{eq:elementary_score}. QLIKE is the Bregman loss of
$\phi(u) = -\log u$, so for every forecast and outcome
\begin{equation}
    \frac{RV}{\widehat{RV}} - \log\frac{RV}{\widehat{RV}} - 1
    \;=\; \int_0^\infty S_\theta\bigl(\widehat{RV},\,RV\bigr)\,
    \theta^{-2}\,d\theta .
    \label{eq:qlike_mixture}
\end{equation}
The identity is a two-line calculation --- when $RV > \widehat{RV}$ the
integral is $\int_{\widehat{RV}}^{RV}(RV-\theta)\,\theta^{-2}\,d\theta =
RV/\widehat{RV} - 1 - \log(RV/\widehat{RV})$, and the other case is its mirror
image --- and it was also checked numerically, day by day, for the
block-diagonal ridge on the scored days. QLIKE is therefore an average of
threshold decisions, weighted by $\theta^{-2}$ toward small thresholds.

The $\mathrm{sign}(s)$ rule takes one of those decisions and no other: it is
long exactly when $\widehat{RV}_t > \mathrm{IV}^2_{30,t}$, the decision at
$\theta = \mathrm{IV}^2_{30,t}$. With the threshold measured in units of the
day's implied variance, $\theta = c\,\mathrm{IV}^2_{30,t}$, the right-hand side
of Equation~\eqref{eq:qlike_mixture} becomes $\int_0^\infty e_t(c)\,c^{-2}\,dc$,
where $e_t(c)$ is the elementary score of $\widehat{RV}_t/\mathrm{IV}^2_{30,t}$
against $RV_t/\mathrm{IV}^2_{30,t}$ at the threshold $c$; the trade reads that
score at $c = 1$ only,
\begin{equation}
    e_t(1) \;=\; \left|\frac{RV_t}{\mathrm{IV}^2_{30,t}} - 1\right|
    \mathbf{1}\bigl\{\widehat{RV}_t \text{ and } RV_t
    \text{ fall on opposite sides of } \mathrm{IV}^2_{30,t}\bigr\}.
    \label{eq:score_at_threshold}
\end{equation}
A single elementary score is itself consistent for the mean, though not
strictly, so it ranks variance forecasts in the same sense QLIKE does; it
spends all of its weight at the threshold the trade uses.
Table~\ref{tab:close_option_scoring} reports it beside QLIKE. QLIKE remains
the paper's measure of accuracy everywhere else.

\begin{table}[H]
\centering
\caption{The 15:30 forecasts scored two ways, on the 866 expiration days of
Table~\ref{tab:close_option_rules}. QLIKE is Equation~\eqref{eq:dp_qlike}
applied to the $16{:}00$ row; the score at the threshold is
Equation~\eqref{eq:score_at_threshold}, averaged over days (lower is better
for both). ``Same side as $RV$'' is the percentage of days on which
$\widehat{RV}_t$ and $RV_t$ fall on the same side of
$\mathrm{IV}^2_{30,t}$; ``same side as $R$'' the percentage on which the
$\mathrm{sign}(s)$ position has the sign of the settlement return. The last
column is the annualized Sharpe ratio of the $\mathrm{sign}(s)$ portfolio at
the crossed spread. The last row is a rule told $RV_t$ in advance. $\dagger$
as in Table~\ref{tab:close_option_rules}.
\label{tab:close_option_scoring}}
\resizebox{\textwidth}{!}{% AUTO-GENERATED by writeup/make_murphy_1530_tex.py — do not edit.
% Source: results/atm_straddle_0dte_1530/proposals/59/a_scores.csv,
%         a_variance_vs_payoff.csv
\begingroup\small\setlength{\tabcolsep}{4.5pt}
\begin{tabular}{lrrrrr}
\toprule
& QLIKE & score at $\theta=\mathrm{IV}^2_{30}$ & same side as $RV$ & same side as $R$ & Sharpe$_{\mathrm{ann}}$, crossed \\
\midrule
baseline (HAR + calendar OLS) & 0.108 & 0.102 & 68.5 & 55.2 & 0.49 \\
block-diagonal ridge & 0.110 & 0.109 & 65.9 & 54.5 & 0.87 \\
block-diagonal ridge, without the FOMC columns & 0.110 & 0.111 & 64.7 & 54.2 & 0.80 \\
LightGBM$^{\dagger}$ & 0.101 & 0.094 & 68.7 & 54.7 & 0.98 \\
XGBoost$^{\dagger}$ & 0.103 & 0.097 & 68.8 & 55.3 & 1.02 \\
lasso (causally tuned)$^{\dagger}$ & 0.104 & 0.106 & 66.3 & 55.1 & 0.93 \\
lasso ($\alpha=10^{-4}$) & 0.111 & 0.107 & 65.7 & 54.0 & 1.04 \\
elastic net (causally tuned)$^{\dagger}$ & 0.104 & 0.109 & 66.5 & 54.6 & 0.54 \\
\midrule
realized variance known in advance & 0 & 0 & 100.0 & 62.9 & 4.07 \\
\bottomrule
\end{tabular}
\endgroup
}
\end{table}

\paragraph{The Murphy diagram.} Averaging $e_t(c)$ over the scored days and
plotting it against $c$, one line per forecast, gives the Murphy diagram of
\citet{Ehm2016} (Figure~\ref{fig:close_option_murphy}). QLIKE is the area
under each line weighted by $c^{-2}$; the trade reads each line at $c = 1$.
The lines cross within a factor of two of the implied variance. At every
threshold on the grid up to \murLinUpper{} times it the lowest score belongs
to a penalized linear forecast --- the block-diagonal ridge, one of the two
lassos or the elastic net. From \murTreeLower{} to \murTreeUpper{} times it
the two tree forecasts lead, \murBestSlice{} at $c = 1$ itself with a score
of \murBestSliceScore{}. Between \murBaseLower{} and \murBaseUpper{} times it
the baseline leads at \murBaseLeadN{} of the \murBaseRegionN{} thresholds, and
from \murTieLower{} times upward the best score is shared exactly by
forecasts of different families, because few forecasts ever reach so high a
threshold and those that do not are scored identically. At the trade's own
threshold the block-diagonal ridge ranks \murRidgeRank{} of the eight, with
\murRidgeSliceScore{}, and the same estimator without the FOMC columns ranks
\murRidgeNoFomcRank{}, with \murRidgeNoFomcScore{}.

\begin{figure}[H]
\centering
\treefigorpending{\textwidth}{figures/murphy_1530.pdf}
\caption{Murphy diagram of the eight 15:30 forecasts on the 866 expiration
days. The horizontal axis is the threshold $\theta$ in units of the day's
implied variance of the last half hour, $\mathrm{IV}^2_{30,t}$, on a
logarithmic scale from \murGridMin{} to \murGridMax{}; the dotted vertical
line at 1 is the threshold the $\mathrm{sign}(s)$ rule uses. (a) The mean of
the elementary score $e_t(c)$ (lower is better); QLIKE is the area under each
line weighted by $c^{-2}$. (b) Each line minus the mean of the eight at the
same threshold, which makes the crossings visible.}
\label{fig:close_option_murphy}
\end{figure}

\paragraph{What a variance forecast can reach.} Even the right threshold
decision does not settle the trade. The straddle's settlement value depends on
where the index closes relative to the two strikes --- on the net move of the
last half hour --- and not on its realized variance, which counts every swing
whether or not it is undone before the close. A rule told $RV_t$ in advance and
trading $\mathrm{sign}(RV_t - \mathrm{IV}^2_{30,t})$ takes the side of the
payoff's sign on only \murOracleAgree{} percent of days, for an annualized
Sharpe ratio of \murOracleShMid{} at the mid and \murOracleShX{} across the
spread, against \murPayoffOracleShX{} across the spread for a rule told the
sign of the payoff itself. On \murOracleMiss{} percent of days --- about a
third --- the trade's sign is therefore beyond any variance forecast, however
accurate. The eight forecasts take the side of the payoff's sign on
\murAgreePayMin{} to \murAgreePayMax{} percent of days, and the side of
$RV_t - \mathrm{IV}^2_{30,t}$ on \murHitVarMin{} to \murHitVarMax{} percent.

\paragraph{Neither score orders the forecasts on this trade.} The fixed lasso
ranks \murLassoFixedQLIKERank{} of the eight on QLIKE and
\murLassoFixedSliceRank{} at the trade's threshold, and has the highest
crossed-spread Sharpe ratio of the eight (\murLassoFixedShX{}). Across the
eight columns the rank correlation between the crossed-spread Sharpe ratio and
minus the loss is \murRhoQLIKE{} for QLIKE (95\% interval from a day-block
bootstrap \murRhoQLIKELo{} to \murRhoQLIKEHi{}) and \murRhoES{} for the score
at the threshold (\murRhoESLo{} to \murRhoESHi{}); both intervals contain zero.
Adding \murNLadder{} deliberately improved and degraded copies of the ridge
forecast --- blends toward the realized value, and multiplicative noise ---
lifts both correlations, to \murRhoUnionQLIKE{} and \murRhoUnionES{}, but that
says only that a much better or much worse forecast trades better or worse:
every loss orders those copies perfectly, so they cannot choose between losses.
The score at the threshold is reported because it isolates the one decision
the trade takes, not because it ranks the eight columns here; with
\murNForecasts{} forecasts that take the same side on most days, and a payoff
whose sign a variance forecast can reach only in part, no variance score does.

\paragraph{The signal's content is a same-day sign, and the right test
is the same-day one.} The signal is fixed at 15:30 and the return
settles at 16:00 on the same expiration day, so pairing $s_t$ with
$R_t$ involves nothing observed after the decision. On that pairing
the content is a shift in the conditional mean by sign: when the
forecast exceeds the implied variance the package returns $+0.10$ on
average, and $-0.09$ when it does not (difference $0.19$,
heteroskedasticity-robust $t=2.32$ for the block-diagonal ridge; $t$
between 1.63 and 2.61 across
the eight columns). That is the $\mathrm{sign}(s)$ portfolio of
Table~\ref{tab:close_option_rules}, read as a regression, and
Figure~\ref{fig:close_option_sameday} shows it. A straight
line is the wrong instrument for the same relation: fitted by least squares
the slope of $R_t$ on the raw signal comes out \emph{negative} for every one
of the eight columns (heteroskedasticity-robust $t$ from $-3.9$ to
$-0.2$) --- the opposite sign to the step the split finds --- because
$R$ has a point mass at $-1$ on 22.3 percent of days and a fat right
tail that a few extreme values of $s_t$ can lever. The effect lives in the
conditional means of the two sides, which a sign split sees and a slope does
not.

\paragraph{Magnitude carries information; sizing on it does not.} Within
either side of the split the binned means of the figure show no trend in
the level of the signal, but a line fitted \emph{inside} a side does find
one, and it runs the wrong way for a trader. Restricted to the sell side
the heteroskedasticity-robust slope of $R_t$ on $s_t$ is negative and
resolved for all eight columns ($|t|$ from 2.2 to 5.7;
$5.4$ for the block-diagonal ridge): the further below zero the
signal, the larger the settlement return, so the short earns \emph{less}
on precisely the days its signal is most extreme. On the buy side the
slope is negative too, but resolved for only two of the eight. The
precise claim is therefore narrower than ``magnitude adds nothing to the
sign'': the conditional mean does vary with $|s|$ inside a side. What does
not follow is a sizing rule, and none is retained --- scaling the position
by the rank of $|s_t|$ raised no column's Sharpe ratio, and that overlay is
parked (Section~\ref{sec:close_option_methods}). Only the sign is used.

\begin{figure}[H]
\centering
\treefigorpending{\textwidth}{../results/atm_straddle_0dte_1530/regression_R_on_signal_blk2.png}
\caption{The 15:30 signal against the 16:00 settlement return, block-diagonal
ridge forecast, 866 expiration days. Left: the settled return $R_t$ against
the raw signal $s_t$ with the least-squares line. Middle: the mean of $R_t$ on
days with $s_t \le 0$ and $s_t > 0$ (means only, no error bars) beside the
always-short portfolio's mean over all days (grey bar), the benchmark throughout.
Right: the mean of $R_t$ in ten equal-count bins of the signal's level, each
labelled by its median value of $s_t$ in units of $10^{-6}$ (thirty-minute
variance); the dashed line is the always-short mean. The relation is a step at
the sign, and the binned means inside either half show no trend in the level:
what the signal carries is which side of the quoted variance the forecast
stands on.}
\label{fig:close_option_sameday}
\end{figure}

\paragraph{At the crossed spread the edge survives and the control does
not.} Every figure above is a midpoint figure. Charging the spread at
entry --- both legs bought at the ask when the rule is long, both sold at
the bid when it is short, with nothing charged at the exit because the
position settles in cash --- takes the block-diagonal ridge
$\mathrm{sign}(s)$ portfolio from an annualized Sharpe of 1.34 to
0.87 ($t=1.61$) and the always-short control from 0.20 to
$-0.27$. That is the comparison the trade turns on: the direction rule
pays for its own spread and the always-short control does not. It is one
crossing per day at the quoted touch, with no size, no queue and no model
of impact --- a bound on the fill rather than a fill model --- but it is
the crossing the venue charges, and the sign of the answer does not depend
on where inside the spread a real fill would land.

\paragraph{The forecast is dated where it claims to be.} A same-day trade
driven by a same-day forecast invites the obvious suspicion, so the shift
test is reported rather than asserted. Re-running the whole portfolio on
the forecast row stamped $k$ bars away from the one the rule uses gives an
annualized Sharpe of 0.48 at $k=-1$ (a forecast issued one bar early,
for a window that has already closed), 1.34 at $k=0$ (the row actually
used), and 3.11 at $k=+1$ --- one bar of genuine look-ahead ---
against a ceiling of 4.53 for a portfolio told the last half hour's realized
variance in
advance. The profile is a step exactly at $k=0$: flat to its left and far
higher to its right. A leak would have lifted the bars left of $k=0$ or
flattened the step, and it does neither. One column is worth flagging: the
elastic net's best stale shift, at $k=-2$, is 1.04 against its own
1.02 at $k=0$, so that one forecast is not separated from its own stale
copies by this test. Every other column's best stale shift sits below its
trade.

% VOL-TARGET PARKED 2026-09-02: overlay held out of the deck and the paper by author order.
\iffalse
\paragraph{Sizing by trailing portfolio volatility standardizes risk without
touching the edge.} Scaling the daily position by the ratio of the
expanding median of the portfolio's own trailing 63-day volatility to its
current lagged value --- self-normalizing, no free target parameter, a
cap of 3 that never binds --- leaves mean and Sharpe statistically
unchanged (block-diagonal ridge $\mathrm{sign}(s)$ 1.94 to 1.98 on the 807
post-warm-up days) while cutting the variability of the portfolio's rolling
volatility by 15 to 21 percent for every forecast and rule and roughly
halving the spread of its per-year volatility (cross-year coefficient
of variation 0.196 to 0.104). This is a standardization of risk, not a
forecast claim, and its structural limit is stated plainly: a trailing
estimator cannot see the first day of a regime change, so single worst
days keep their size (excess kurtosis rises from 10.6 to 12.8) and the
COVID transition itself falls inside the estimator's warm-up.
\fi

% IRON-FLIES PARKED 2026-09-03: cost at every width, fraction of wealth unchanged; deck section parked, lab remains in the experimental notebook.
\iffalse
\paragraph{The price of the defined-risk constraint.} Per dollar of
body premium the wings cost the $\mathrm{sign}(s)$ portfolio at every width, and
the cost falls with distance: annualized Sharpe 1.36, 1.42, 1.47 and
1.54 at $w = 20, 25, 30, 50$ against 1.62 for the plain portfolio on the
same 870 days (paired daily difference $-221$, $-164$, $-120$ and $-62$
basis points; Newey--West $t$ between $-2.4$ and $-3.0$ for the
block-diagonal ridge, and a lower Sharpe ratio for every one of the seven
forecasts at every width). The always-short control does not survive the constraint at
any width below 50 ($-0.15$, $-0.06$, $0.00$ and $+0.10$ against
$+0.27$). What the constraint buys is the dollar tail on the days the
$\mathrm{sign}(s)$ portfolio sells:
the worst settlement falls from $-78$ points per package to $-17$,
$-20$, $-25$ and $-37$, with the settlement landing beyond a wing on
7.3, 4.3, 2.5 and 1.1 percent of selling days. Per premium the
$\mathrm{sign}(s)$ portfolio's worst day barely moves ($-4.7$ at $w=20$; unchanged
at $-5.4$ for any wider wing), because its worst days are moderate
moves on low-premium days that no wing reaches. At midpoint quotes the
wings roughly pay for themselves in points --- premium paid against
settlement received on selling days is within a few points at every
width, $t$ near zero --- but that is one quarter's doing: the first
quarter of 2020, inside the estimation warm-up, returned 50 to 83
percent of everything the wings ever paid back, and in 2021 and 2024
they paid nothing at any width of 25 or more. At realistic fills ---
the body sold at the bid, the wings bought at the ask, whose median
sits 20 to 100 percent above the midpoint --- the price is higher: the
$\mathrm{sign}(s)$ portfolio's crossed-spread Sharpe is 0.69, 0.77, 0.83 and 0.92 against
1.11 for the plain portfolio at the crossed spread.
\fi

\paragraph{Compounded at a fixed three percent of wealth, the $\mathrm{sign}(s)$
portfolio grows; the control does not.} Betting the fixed fraction
$f=0.03$ of Section~\ref{sec:close_option_methods}, the block-diagonal ridge
$\mathrm{sign}(s)$ portfolio turns one unit of wealth into $7.2$ over the
sample (annualized log-growth $+0.57$), with a maximum wealth drawdown
of 42 percent and a worst single day that removes 16 percent of
wealth; the ruin bound $1/|\min_t R'_t|$ is 0.18 for this portfolio and
0.10 for the control, so three percent sits well inside both. All eight
$\mathrm{sign}(s)$ portfolios compound positively under the same fraction (terminal
wealth $3.6$ to $9.9$), and the headline portfolio's annualized
log-growth is positive in four of the five calendar years (2022 is $-0.04$). The
always-short control ends at $0.86$ --- its premium does not compound. A defined-risk variant is
parked and not reported here.

% IRON-FLIES PARKED 2026-09-03: cost at every width, fraction of wealth unchanged; deck section parked, lab remains in the experimental notebook.
\iffalse
\paragraph{Under the constraint the fraction of wealth barely moves,
and the frame decides the number.} The bounded worst day is exactly
what would loosen the ruin bound on the fraction, but in this sample
the bound was never the binding term: the causal fraction of the parked estimated-fraction rule sits well below it
on every day for
the plain portfolio (mean 0.063, maximum 0.106) and stays there under the
wings (mean 0.054 to 0.062 per premium across the four widths; the
running cap binds on no day in any portfolio). The fraction is set by the
estimate of mean over second moment, not by the worst day, so capping
the worst day does not let the trader bet more. In the per-premium
frame the defined-risk portfolio compounds to 10.3, 12.3, 14.8 and 21.9
times initial wealth at $w = 20, 25, 30, 50$ against 33.9 for the plain
portfolio on the same days --- the wing cost compounds into a real growth
cost, smallest at the widest wing --- while the worst single-day wealth
factor improves only from 0.51 to between 0.53 and 0.64. In the
capital-at-risk frame, well defined here because no fraction below one
can be ruined, the same portfolio is a fully collateralized floor: a mean
fraction of 0.03 to 0.04 of wealth posted as collateral, terminal
wealth 2.3 to 2.7, worst single day $-7$ percent of wealth. The two
frames are different units and compare only within themselves. Taken
together the numbers support the widest wing the margin rule allows:
at $w=50$ the constraint costs 0.08 of Sharpe and a third of compounded
growth while halving the worst selling day in points, and every
narrower wing buys a smaller dollar tail at two to three times that
cost. The width is a choice of the dollar worst day the account can
carry, not a growth decision --- and every figure here is a midpoint
figure that worsens at the ask.
\fi

The trade as a whole is a demonstration under the stated constants of
Section~\ref{sec:close_option_methods}, not a trading claim: there is
no delta hedge, no borrow and no market-impact model, and the headline
fills are at the 15:30 mid. What survives the qualifications is worth
stating plainly. At the last bar the quoted implied variance sits above the
forecast on 59 to 68 percent of days depending on the column
(60 for the headline); the sign of that gap carries a direction
edge, one that improves on the always-short control in every one of the
eight columns and survives a crossed spread --- though only one column's
improvement clears zero on a bootstrap interval, and that by a knife edge,
and the columns cannot be ordered against one another; and sizing on the
magnitude of the gap does not improve on its sign.
